% !TeX encoding = utf8

\documentclass[a4paper,8pt]{article} %en vez de 11pt
\usepackage[spanish]{babel}
%\usepackage[english]{babel}
\usepackage[utf8]{inputenc}
\usepackage{amsmath,amssymb,amsthm,xspace,a4wide}
\usepackage{hyperref,color}
\usepackage[a4paper, total={6in, 8in}]{geometry}
%\usepackage{macros}
\usepackage{enumitem}% Permite poner cosas del tipo \begin{enumerate}[label=\Alph*]
\usepackage[table]{xcolor}
\usepackage{colortbl}


%\usepackage{pgfgantt}
%\usetikzlibrary{decorations.pathmorphing, shadows, shadings}

\hypersetup{
    colorlinks=true,   % color instead of boxes
    citecolor=blue,    % cite links in blue
    linkcolor=blue,    % other internal links in gray
%    linktocpage,         % in TOC, LOF and LOT, the link is on the page number
}


\usepackage{sectsty}%Para cambiar tamaño de Section sizes
\sectionfont{\fontsize{8}{15}\selectfont}



\usepackage[font=small,labelfont=bf]{caption} % Para el texto de las figuras sea chico es el font = small. Para que Figure este en mayuscula seria con labelfont = bf


\let\OLDthebibliography\thebibliography %para reducir espacio en biblio
\renewcommand\thebibliography[1]{
  \OLDthebibliography{#1}
  \setlength{\parskip}{0pt}
  \setlength{\itemsep}{0pt plus 0.3ex}
}


\begin{document}
\thispagestyle{empty} %Para no poner page number

\begin{center}
% {\huge Entendiendo los límites teóricos de los LLMs desde el punto de vista de la complejidad computacional: el rol de la aleatoriedad dentro del Chain Of Thought}

% \vspace{0.5cm}
{\huge \textbf{Plan de trabajo}}

\vspace{0.2cm}
{\large Tesis de Licenciatura en Ciencias de la Computación}
\end{center}
\vspace{-0.1em}
\noindent
\vspace{-0.1em}\textbf{Estudiante}: Rafael Romani \quad \texttt{rafaromani243@gmail.com} \quad LU 775/21 \\
\vspace{-0.1em}\textbf{Director}: Santiago Cifuentes \quad \texttt{santiagocifuentes66@hotmail.com} \\
\vspace{-0.1em}\textbf{Codirector}: Santiago Figueira \quad \texttt{sfigueir@gmail.com}

\section*{\Large Título tentativo}
\noindent “Entendiendo los límites teóricos de los LLMs desde el punto de vista de la complejidad computacional: el rol de la aleatoriedad dentro del Chain Of Thought"

\section*{\Large Contexto}

En los últimos años, los modelos de procesamiento del lenguaje natural y las inteligencias artificiales generativas comenzaron a cobrar un rol central en la sociedad.
La mayoría de estos modelos están construidos sobre \textit{transformers}, una arquitectura para modelos de NLP presentada en 2017 por investigadores de Google \cite{ATTN-all-you-need}.
Un \textit{transformer} consiste en múltiples capas en las que se mapea una secuencia de \textit{tokens} a vectores sobre los cuales se computan diversos algoritmos para luego definir un \textit{token} de salida.

Una técnica conocida para mejorar la performance de los grandes modelos del lenguaje (LLM) es forzarlos a imprimir pasos intermedios en su razonamiento.
Este proceso se conoce como \textit{Chain of Thought} (CoT).
Se puede hacer que los modelos impriman pasos intermedios sin necesidad de entrenarlos nuevamente. 
Para conseguirlo, basta con agregar frases a la consulta como \textit{let's think step by step} o incluso agregar los primeros pasos del razonamiento (aunque sean equivocados) \cite{wang-etal-2023-towards} \cite{cot-prompt}.
Esto es conocido como \textit{CoT-prompting}.

\medskip

Observando los incrementos agigantados en las capacidades de los LLM, naturalmente surge la pregunta y preocupación sobre cuál es el límite de su poder y qué se requiere para alcanzarlo.
Múltiples estudios intentaron responder de manera teórica a estas preguntas, pero la representación de los \textit{transformers} como máquinas formales es un desafío aún abierto.
En la literatura se encuentran diversas opciones.
Cada formalización se toma distintas libertades para encontrar un punto medio entre representar fielmente a un \textit{transformer} real y definir una máquina lo suficientemente sencilla sobre la que se puedan probar teoremas.
Entre las diferencias que se observan aparecen cuestiones como qué limitaciones tienen las funciones utilizadas para transformar las cadenas de \textit{tokens} en cadenas de vectores o qué sistema se utiliza para representar los números que componen a los vectores.
Según estas decisiones, se obtienen resultados distintos para el poder expresivo de los \textit{transformers}.

\medskip

El poder de las máquinas de Turing se analiza en relación con cuántos pasos se las deja computar o cuánta memoria utilizar.
En los \textit{transformers} no hay consenso sobre cuál es el recurso a estudiar.
Por ejemplo, Merrill y Sabharwal analizaron qué ocurre al utilizar modelos con más capas \cite{log-depth}.
En cambio, Li et al estudiaron qué clases de problemas pueden resolver cuando permiten más o menos pasos de CoT pero dejando fija la cantidad de capas. Además, analizaron qué ocurre al permitir agrandar la dimensión de los vectores con los que los modelos trabajan \cite{toyota}.


Dadas las diferencias entre formalizaciones y poco consenso en qué recurso limitar para estudiar el poder de cómputo, adaptar resultados entre distintas investigaciones se vuelve complejo.
El área no tiene definiciones claras ni simulaciones entre las distintas opciones.
Además, las formalizaciones estudiadas son todas determinísticas.
Por lo tanto, no capturan el comportamiento estocástico que intrínsecamente tienen los \textit{transformers}.




\section*{\Large Objetivos y plan}

Esquemáticamente, se espera llevar a cabo el siguiente plan:

\begin{enumerate}
    \item Realizar una revisión bibliográfica para identificar las diversas formalizaciones que se utilizan y los resultados de expresividad asociados a ellas.
    \item Unificar las formalizaciones que sea posible.
    \item Estudiar las propiedades de las clases de problemas que se pueden resolver al limitar los pasos de CoT.
    \item Definir una formalización probabilística para los \textit{transformers}.
    \item Definir clases de problemas que los \textit{transformers} probabilisticos pueden resolver según los pasos de CoT que se les permite hacer.
    \item Estudiar las propiedades de estas clases y analizar si coinciden con clases conocidas.
\end{enumerate}

\section*{\Large Cronograma y Factibilidad}

El tesista ya viene trabajando en el plan delineado en la sección anterior a través de una beca BIIC - Lambda Class desde septiembre de 2025. 
Dicho eso, las únicas tareas restantes son aquellas relacionadas con estudiar las clases de complejidad probabilísticas.
Se espera que estas tareas, junto a la escritura de la tesis, le demanden al tesista 3 meses.



\vspace{1em}

%\vspace{-1em}
\renewcommand{\section}[2]{}% Para que no ponga el título Referencias
%\sectionfont{\fontsize{10}{15}\selectfont}

{\small 
\bibliographystyle{plain-initials} % Para que ponga inicial de nombre + apellido de forma estandarizada
\bibliography{biblio}
\thispagestyle{empty}
}


\end{document}
